
\section{Conclusion}
\label{sec:conclusion}
We measured how consumer IoT devices receive software updates in our own update experiments and in an earlier public dataset, and separated update traffic from everything else before assessing it.
This separation matters: keyword matching selects protocol vocabulary rather than updates, and statistics over whole captures mix update channels with video streams, local control traffic, and the certificates of unrelated services.
On the update channels themselves, three of the seven devices that downloaded an update received the image unencrypted\textcolor{blue}{, and unencrypted update requests disclose model and firmware version, in 2019 together with MAC addresses}.
Encrypted channels mostly use strong settings, with outdated key exchange and cipher modes as exceptions, and the legacy options that devices offer mirror software defaults rather than what is actually used.
Weak or long-lived server certificates appear only where vendors run their own certificate authorities\textcolor{blue}{, and the devices that trust them never checked whether a certificate had been revoked}.
Whether these weaknesses let an attacker alter an update depends on the signature inside the image, which the two unencrypted downloads carry, and the local push appears to lack, so the channel and the image have to be assessed together.

% ---- Acknowledgments ----
%\begin{credits}
%\subsubsection{\ackname} A bold run-in heading in small font size at the end of the paper is
%used for general acknowledgments, for example: This study was funded
%by X (grant number Y).
%\end{credits}


%Against the 2019 traces this posture has \emph{moved} rather than improved outright: TLS~1.3 now appears where it did not.
%}
%, while the cheapest cameras remain without forward secrecy.
%That the remedies are configuration-level, dropping legacy suite advertisement, shortening certificate lifetimes, and encrypting firmware transfers, makes the gap a matter of vendor practice rather than device capability.}


%An adversary who can observe update traffic can fingerprint devices, extract weak cryptographic parameters, and prioritize targets for further exploitation, a capability that persists even when on-device firmware signature verification is in place.



% Well, i just did the below limitation

%Third, perform controlled, instrumented experiments in laboratory settings.  Run actual update campaigns on instrumented devices to observe the network traffic negotiation outcomes, actual artifact retrievals, signature verification, and certificate-chain validation,  and on-device installation behavior. These experiments will (i) validate entropy-based thresholds, (ii) confirm whether weak cipher offers are actually negotiated in practice, and (iii) enable end-to-end assessment of update integrity and resilience to active attacks.

%\clearpage
%\appendix
%\input{Appendix}




% ---- Acknowledgments ----

%\begin{credits}
%\subsubsection{\ackname} A bold run-in heading in small font size at the end of the paper is
%used for general acknowledgments, for example: This study was funded
%by X (grant number Y).

%\end{credits}

